\begin{filecontents}[overwrite]{compiler.bib}
@book{appel1998tiger,
	author    = {Andrew W. Appel},
	title     = {Modern Compiler Implementation in ML},
	publisher = {Cambridge University Press},
	year      = {1998},
	isbn      = {978-0521607643},
	note      = {Tiger book}
}
@book{siek2023essentials,
	author    = {Jeremy G. Siek},
	title     = {Essentials of Compilation: An Incremental Approach in Racket},
	publisher = {The MIT Press},
	year      = {2023},
	isbn      = {978-0262047760},
	address   = {Cambridge, Massachusetts}
}
@book{appel1992compiling,
	author    = {Andrew W. Appel},
	title     = {Compiling with Continuations},
	publisher = {Cambridge University Press},
	year      = {1992},
	address   = {Cambridge, England},
	isbn      = {0-521-41695-7}
}
@online{wiki_tail_call,
	title   = {Tail call},
	url     = {https://en.wikipedia.org/wiki/Tail_call},
	urldate = {2026-10-14}
}

@online{wiki_cps,
	title   = {Continuation-passing style},
	url     = {https://en.wikipedia.org/wiki/Continuation-passing_style},
	urldate = {2026-10-14}
}
\end{filecontents}

\documentclass{article}
\usepackage{fontspec}%[no-math]
\usepackage{blindtext}
\usepackage{polyglossia}
\usepackage{hyperref}

\usepackage{listings}
\usepackage{caption}
\usepackage[a4paper, left=3cm, top=1.5cm, right=3cm]{geometry}

\usepackage[style=authoryear, dashed=false]{biblatex}
\addbibresource{compiler.bib}

\setdefaultlanguage{russian}
\setotherlanguage{english}

\setmainfont{Source Sans 3}
\setsansfont{Source Sans 3}
\setmonofont[Scale=0.9]{Roboto Mono}

\renewcommand{\lstlistingname}{Пример}
\captionsetup[lstlisting]{
	justification=centering,
	singlelinecheck=false
}

\title{Оптимизация хвостовой рекурсии в Rukaml.\\Документация}
\author{Егоров Никита}
\date{}
\begin{document}
	
	\maketitle
	\tableofcontents
	
	\section{Введение}
	% В начале я бы не хотел обращать внимание читателя на проблему, где лежат аргументы функции, поскольку это требует глубокого рассмотрения.
	
	% Чтобы понимать, как устроена оптимизация, неплохо не только понимать как её использовать, но знать как она выглядит на уровне ассемблера.

	Хвостовой вызов — это вызов функции, после которого вызывающей функции больше ничего не нужно делать.
	
	\noindent
		\begin{minipage}[t]{0.48\textwidth}
		\begin{lstlisting}[
			language=ml,
			captionpos=b,
			label={lst:1},
			caption={Не хвостовой вызов}]
let rec fac n =
  if n <= 1 then 1
  else n * fac (n - 1)
		\end{lstlisting}
	\end{minipage}
	\hfill
	\begin{minipage}[t]{0.48\textwidth}
		\begin{lstlisting}[
			language=ml,
			captionpos=b,
			label={lst:2},
			caption={Хвостовой вызов}]
let rec fac n acc =
  if n <= 1 then acc
  else fac (n - 1) (n * acc)
		\end{lstlisting}
	\end{minipage}

	Далее \texttt{caller} - вызывающая функция, \texttt{callee} - вызываемая функция.
	
	Хвостовой вызов позволяет избежать аллоцирования стекового кадра \texttt{callee}, переиспользовав кадр \texttt{caller}-а. Таким образом, во время исполнения \hyperref[lst:2]{пример №2} будет занимать лишь $O(1)$ памяти, когда как \hyperref[lst:1]{пример №1} $-$ $O(n)$.
	
	Примечание: хвостовой вызов, например, инструкция \texttt{tail} в ассемблере \texttt{RV64} $-$ это прыжок без возврата, когда как обычный вызов, например, \texttt{call} - это прыжок с возвратом. Важно понимать эту разницу.
	

\section{Как её использовать}

Приведённый \hyperref[lst:2]{пример №2} является демонстрацией работы оптимизации для функций верхнего уровня.

В модуле 2.2 будет рассмотрен следующий по частоте вид - передача функции как аргумента.

А в модуле 2.3 будет показан способ, позволяющий программно преобразовать любую функцию так, чтобы в каждой ветви исполнения она использовала хвостовой вызов.

\subsection{Функции верхнего уровня}

Покажем дополнительный пример.

\begin{lstlisting}[language=ml]
let rec foldr f ini = function
| [] -> ini
| x :: xs -> f x (foldr f ini xs)

let rec foldl f ini = function
| [] -> ini
| x :: xs -> foldl f (f ini x) xs
\end{lstlisting}

Два приведённых выше примера - правая и левая свёртка соответственно.

\texttt{foldr} не вызывается хвостовым образом, поскольку после возврата из функции полученное значение подставляется в \lstinline|f x ...|.

С \texttt{foldl} проблем никаких нет.
\subsection{Замыкания}

Существуют случаи, когда нам необходимо вызвать переданную как аргумент функцию.
\begin{lstlisting}[
	language=ml,
	captionpos=b,
	label={lst:3},
	caption={Факториал в CPS стиле}]
let rec fac_cps n k = 
  if n <= 1 then k 1
  else fac_cps (n - 1) (fun r -> k (n * r))
\end{lstlisting}

Существенная разница между 2.1 и 2.2 появляется, когда дело доходит для реализации, что будет описано позже.

\subsection{CPS}
Continuation-passing style — стиль программирования, в котором поток управления передаётся в явном виде через вызов “продолжения”.

CPS функция принимает дополнительный аргумент: явное «продолжение» — функцию одного аргумента, которая будет вызвана по окончании вычислений с их результатом в качестве аргумента.

Замечание: каждый вызов в CPS является хвостовым.

Пусть у нас есть следующие функции:
\begin{lstlisting}[language=ml]
let pow2 a = a ** 2.0
let add a b = a +. b
let sqrt a = a ** 0.5
let hyp a b = sqrt (add (pow2 a) (pow2 b))
\end{lstlisting}

Перепишем их в CPS-стиле:
\begin{lstlisting}[language=ml]
let pow2 a k = k (a ** 2.)
let add a b k = k (a +. b)
let sqrt a k = k (a ** 0.5)
let hyp a b k = 
  pow2 a (fun a2 -> 
    pow2 b (fun b2 -> 
      add a2 b2 (fun c2 -> sqrt c2 k)))
\end{lstlisting}

Примечание: преобразовать мы можем любую функцию, в том числе и которая уже была преобразована.
Например, преобразованный компилятором Rukaml  \hyperref[lst:3]{пример №3}:
\begin{lstlisting}[language=ml]
let rec fac_cps n k1 = 
  k1 (fun k -> 
    fun k2 -> if n <= 1 then k 1 k2 
              else 
              fac_cps (n - 1) 
                      (fun t3 -> t3 (fun r -> fun k4 -> k (n * r) k4) k2)
     )
\end{lstlisting}

\section{Как она реализуется в Rukaml}

На данный момент в backend действуют разные соглашения о вызовах.

\begin{itemize}
	\item RV64 - первые 7 аргументов через регистры, остальные через кучу.
	\item AMD64 - все аргументы передаются через стек (кроме runtime-функций).
\end{itemize} 

Оптимизация в каждом backend реализована по-разному, но общее в любой реализации:

\begin{enumerate}
	\item Вычислить и переместить аргументы на нужное место.
	\begin{itemize}
		\item AMD64: перемещаем аргументы \texttt{callee} туда, где лежат аргументы \texttt{callee}
		\item RV64: перемещаем первые 7 аргументов через \texttt{a0..a6}, остальные как \texttt{a7[i-8]}, где \texttt{i} $-$ индекс текущего аргумента.
	\end{itemize}
	\item Восстановить \texttt{callee-saved} регистры, \texttt{sp} и адрес возврата.
	\item Прыгнуть без возврата в \texttt{callee}.
\end{enumerate}
\subsection{RV64}

В RV64 всё весьма относительно тривиально:
\begin{itemize}
	\item Для \href{https://github.com/KrimPast/rukaml/blob/7d05249ecd7c3cf5490ba5d807674cd9fc18ba7b/back/rv64/RV64_impl.ml#L1054-L1063}{функций верхнего уровня} $-$ вычисляем аргументы \texttt{callee}, грузим \texttt{callee-saved} аргументы(\texttt{s0..s7}), \texttt{ra} и возвращаем \texttt{sp} как он был на входе в функцию.
	\item Для \href{https://github.com/KrimPast/rukaml/blob/7d05249ecd7c3cf5490ba5d807674cd9fc18ba7b/back/rv64/RV64_impl.ml#L1099-L1104}{замыканий} всё аналогично.
	
Примечание: мы грузим \texttt{callee-saved} регистры только перед хвостовым вызовом и в конце функции.
\end{itemize} 

\subsection{AMD64}
\label{subsec:AMD64_ASM}

В AMD64 всё куда интереснее, поскольку там всё передаётся через стек.
Это делает оптимизацию хвостового вызова на порядок сложнее. Разберём почему это так.

\begin{itemize}
	\item Пусть $ sp(f) = $ количеству аргументов функции f. Все аргументы занимают по машинному слову, так что нет смысла оперировать их размером.
	\item Если $sp(f) \geq sp(g)$, то хвостовой вызов возможен $-$ у нас есть достаточно места, куда сложить аргументы $g$.
	\item Если $sp(f) < sp(g)$, то у нас есть выбор.
	\begin{enumerate}
		\item Отказаться от хвостового вызова, ведь $sp(g) - sp(f)$ аргументов ложить некуда.
		\item Сделать так, чтобы место было.
	\end{enumerate}
	
	Очевидно, был принят второй вариант. 
\end{itemize}

Для реализации этого, составляется табличка "кому сколько нужно дополнительного места".
Она создаёт путём итерирования по всем \texttt{caller} и их хвостовых \texttt{callee}.
Например, для каждой такой пары устанавливается $sp(caller) = max(sp(caller),\ sp(callee))$\footnote{Реализацию таблички можно посмотреть в файле \href{https://github.com/KrimPast/rukaml/blob/7d05249ecd7c3cf5490ba5d807674cd9fc18ba7b/middle/opts.ml}{middle/opts.ml}, в модуле TailCall.}

Кажется всё прекрасно, но у такого способа есть свои грабли.
Когда дело доходит до замыканий, всё ломается.

\begin{itemize}
	\item Пусть функции $f$ передали в качестве аргумента функцию $g$.
	Мы хотим вызвать хвостовым образом функцию $g$, передав туда $n$ аргументов.
	\item Проблема в том, что $g$ $-$ замыкание. У него $m$ аргументов уже применено.
	Однако на стадии компиляции число $m$ в общем случае неизвестно. \footnote{Иначе, чтобы составить такую табличку, нужно пробежаться по абсолютно всему коду и смотреть, скольки аргументное замыкание мы передаём данной функции. Я не утверждаю, что теоретически это невозможно, скорее, если это и выполнимо, то очень трудоёмко.}
\end{itemize}

Таким образом, данный способ в полной мере работает только для функций верхнего уровня.

В реализации для замыканий динамически проверяется, есть ли в текущем кадре достаточно места, и если да, то вызывается хвостовым образом.

\subsection{Почему в RV64 всё так просто}
На первый взгляд кажется, что особо нет разницы $-$ передавать аргументы через стек или через кучу.

Однако, если мы передаём аргументы через кучу, то у нас пропадает проблема, описанная в \ref{subsec:AMD64_ASM}. Поскольку 8-й и следующие хранятся в списке на куче, ссылка на который лежит регистре \texttt{a7}. Например, в общих чертах так и реализовано соглашение о вызовах в \texttt{Scheme/Racket}.
\section{Future Work}
\subsection{Оптимизация хвостового вызова}
\begin{itemize}

\item AMD64: Если будет соглашение о вызовах как сейчас в RV64, то нужно добавить полную поддержку хвостового вызова замыканий.
\item RV64, AMD64: В результате применения оптимизации в коде появляется дополнительное восстановление \texttt{ra}, \texttt{sp }и \texttt{s0..s7} регистров. То же самое происходит и в конце функции перед её возвратом. Можно попробовать избежать этого дублирования.
\end{itemize}

\subsection{Изменение соглашения о вызовах}
\begin{itemize}
	\item RV64: В целях оптимизации производительности переписать \texttt{apply1}, ... , \texttt{apply5} так, чтобы они не использовали applyN.
	\item RV64: Удалить в комментариях упоминания о передаче данных через стек и добавить о регистрах и куче. Ибо это может сбить с толку.

\urldef{\urlsimlify}\url{https://github.com/KrimPast/rukaml/blob/7d05249ecd7c3cf5490ba5d807674cd9fc18ba7b/back/rv64/RV64_impl.ml#L836}

	\item RV64: В коде появились участки, которые можно упростить благодаря новому соглашению о вызовах.\footnote{\urlsimlify}
	\item Переписать(?) Calling Convention под AMD64, чтобы оно тоже использовало регистры и кучу. Эта работа не очень пыльная, если сверяться с тем, как реализовано в RV64.
\end{itemize}

\subsection{Иное}
\begin{itemize}
	
\urldef{\urlargonerv}\url{https://github.com/KrimPast/rukaml/blob/7d05249ecd7c3cf5490ba5d807674cd9fc18ba7b/back/rv64/RV64_impl.ml#L1094}
\urldef{\urlargoneamd}\url{https://github.com/KrimPast/rukaml/blob/7d05249ecd7c3cf5490ba5d807674cd9fc18ba7b/back/amd64/amd64_impl.ml#L906-L908}
	\item Переписать в backend AMD64, чтобы вместо \texttt{"printfn ppf ..."} код использовал \texttt{emit} прямо как в RV64.
	\item AMD64, RV64: В backend применение замыканий пользовательских функций является одноаргументным \footnote{\urlargonerv} \footnote{\urlargoneamd}. Неудивительно, что это ест производительность.
\end{itemize}

\defbibenvironment{itemize}
{\begin{itemize}}
	{\end{itemize}}
{\item}

\nocite{*}
\printbibliography[heading=bibnumbered, title={Полезные источники}, env=itemize]
\end{document}